\section{时间推进与推理}

\subsection{在线查询的分层顺序}
对一个静态 geometry $g$，在线推理先确定性构造 geometry-aware thermal objects：
\begin{equation}
\Phi(g),\qquad M_r(g),\qquad K_r(g),\qquad \phi_j^{\min/\max}(g),
\end{equation}
同时由 neural surrogate 一次得到
\begin{equation}
\{\widehat Z_{\rm field}(g),\widehat D_{\rm vol}(g),\widehat H_1(g),\ldots,\widehat H_r(g)\}.
\end{equation}
两条 geometry-dependent map 在物理上并行：前者是确定性 thermal coordinate construction，后者是学习到的 EM tensor map。

当前 constitutive model 下，上述对象在同一个静态 geometry 的 thermal trajectory 中保持不变，因此全部只需构造一次并缓存。

\subsection{初始条件接口}
若输入 full temperature-rise field $\theta_0$，默认按当前 geometry 的 thermal mass inner product 投影：
\begin{equation}
\boxed{
a_0=[\Phi(g)^TM_T(g)\Phi(g)]^{-1}\Phi(g)^TM_T(g)\theta_0.}
\end{equation}
并报告 initial projection error。若直接输入 reduced coordinates $a_0$，只保证 reduced state 定义明确；只有属于冻结/审计 initial-condition family 才携带 full-order accuracy 声明。

\subsection{规定电流工况}
对给定端口峰值复电流 $c$，体 Joule reduced source 为
\begin{equation}
\widehat q_{{\rm vol},j}(g,c)
=\frac12\operatorname{Re}(c^H\widehat H_j(g)c),
\end{equation}
体总耗散为
\begin{equation}
\widehat P_{\rm vol}(g,c)=\frac12c^H\widehat D_{\rm vol}(g)c.
\end{equation}
thermal ODE 为
\begin{equation}
\boxed{
M_r(g)\dot a
=-K_r(g)a
+\widehat q_{\rm vol}(g,c)
+q_{\rm wire}(a,g,c)+f_T(g).
}
\end{equation}
若 $c=c(t)$ 在 thermal time scale 上缓慢变化，则在每个 integration stage 使用当前 $c(t)$。

\subsection{电压/负载控制必须显式接 circuit equations}
若实际工作点由 source voltage、matching/compensation network 与 load 决定，则端口电流由确定性 circuit equations 求得，例如
\begin{equation}
\boxed{
[\widehat Z_{\rm field}(g)+R_{\rm wire}(a,g)+Z_{\rm ext}(\omega)]c
=v_{\rm src}.
}
\end{equation}
此时 $c=c(a,g)$，thermal dynamics 成为闭环非线性系统。circuit physics 不由网络重新学习。

\subsection{wire temperature 与重构}
对第 $p$ 个线圈，temperature state 通过
\begin{equation}
\bar\theta_p(a,g)
=\ell_p^{(T)T}(g)\Phi(g)a
\end{equation}
进入 $R_p$。full temperature reconstruction 为
\begin{equation}
\boxed{T(x,t)=T_{\rm amb}+\Phi(x;g)a(t).}
\end{equation}
因此 local modes 随 TX/RX pose 变化后，热点空间位置会自动跟随 geometry，而不是依赖 fixed global basis 记忆所有位置。

\subsection{ETD2、IMEX 与固定 geometry}
对静态 geometry，$M_r(g),K_r(g)$ 在时间上固定。定义
\begin{equation}
L(g)=M_r(g)^{-1}K_r(g).
\end{equation}
中低 rank 下可对 generalized eigenproblem
\begin{equation}
K_rv=\lambda M_rv
\end{equation}
做一次分解并复用于整个 trajectory。规定常电流时可使用 ETD2 predictor-corrector；对 $c(t)$ 或 circuit-controlled mode，stage 必须重新评估 current 与 source。若非线性/非自治程度较强，可采用 linearly implicit IMEX 或高精度 stiff reference solver。

\subsection{长时间查询使用连续/链式推进}
finite-time 查询 horizon 不决定 thermal basis dimension。只要 geometry-aware ROM 已在可解析短/中时间尺度和 steady limit 上通过 held-out trajectory audit，就可用同一个稳定 ODE 连续推进到更长时间：
\begin{equation}
a(t_0)\rightarrow a(t_1)\rightarrow\cdots\rightarrow a(t_N).
\end{equation}
因此查询 $1000\,$s、$10000\,$s 不要求专门增加对应的 resolvent anchor。

这里的“链式”是对同一个 reduced ODE 做数值时间推进，不是反复调用 neural network。对固定 geometry，EM tensors 与 thermal reduced operators 都缓存不变。

\subsection{自适应步长}
自适应 ETD/IMEX 只控制时间离散误差。对特别长时间查询，接近稳态后允许步长增大，但必须持续检查 state-domain、circuit solvability、wire constitutive validity 和数值稳定性。时间积分误差必须与 tensor surrogate error、thermal ROM error 分开报告。

\subsection{生产时间积分器的独立 reference audit}
生产 release 不能只验证 reduced vector field 正确，而不验证实际在线 integrator。对 completely-held-out geometry 和冻结的 current-controlled/circuit-controlled operating cases，先固定同一组 surrogate tensors、$\Phi(g)$、$M_r(g)$、$K_r(g)$ 与初值，再分别用生产 integrator 和高精度 stiff reference solver 推进同一个 reduced ODE。

记 reference state 为 $a_{\rm ref}(t)$、生产积分结果为 $a_{\rm prod}(t)$，至少使用当前 reduced thermal mass 定义
\begin{equation}
E_{\rm time}(t)
=\frac{\|a_{\rm prod}(t)-a_{\rm ref}(t)\|_{M_r}}
{\max(\|a_{\rm ref}(t)\|_{M_r},\varepsilon_{\rm abs})}.
\end{equation}
该误差只属于 time discretization，不得混入 tensor surrogate error 或 full-vs-ROM thermal error。生产 artifact 发布前必须满足冻结的 integrator tolerance；失败时拒绝发布，而不是通过放宽 neural/ROM 指标掩盖时间推进误差。

\subsection{稳态查询}
$t=\infty$ 不通过超长时间积分逼近。current-controlled equilibrium 满足
\begin{equation}
-K_r(g)a_*+\widehat q_{\rm vol}(g,c)+q_{\rm wire}(a_*,g,c)+f_T(g)=0.
\end{equation}
circuit-controlled mode 需与 circuit equations 联立求解 $(a_*,c_*)$。

候选 equilibrium 必须满足 residual、constitutive domain 与 local stability。闭环 Jacobian 定义为
\begin{equation}
J_*=M_r^{-1}\frac{\partial}{\partial a}
[-K_ra+q_{\rm vol}+q_{\rm wire}+f_T]_{a=a_*}.
\end{equation}
只有
\begin{equation}
\max\operatorname{Re}\lambda(J_*)<0
\end{equation}
到规定 margin 时才标记为局部渐近稳定 steady state。多个 stable equilibria、thermal runaway 或 no equilibrium 必须显式报告。

\subsection{批量几何查询}
对同一 geometry 的多个时间、电流、外部电路或初值查询，$\Phi(g),M_r(g),K_r(g),\widehat Z,\widehat D,\widehat H_j$ 全部只构造一次。对批量不同 geometry，可批量执行 MLP，同时独立确定性构造各自 thermal basis/ROM。

在线主要成本因此是一次 geometry surrogate forward、一次小规模 thermal projection/assembly 和后续 small thermal/circuit operations，而不是 online Maxwell solve。

\subsection{运动 geometry 不属于当前静态 ROM}
若未来 $g=g(t)$，则
\begin{equation}
\theta=\Phi(g(t))a(t)
\end{equation}
给出
\begin{equation}
M_r\dot a
+\underbrace{\Phi^TM_T\dot\Phi}_{G_r(g,\dot g)}a
+K_ra=q.
\end{equation}
同时 EM tensors、wire functionals 与 reduced operators 也随时间更新。因此运动 geometry 需要显式 transport term 与相应时间尺度验证，不能在当前静态 query engine 中简单每步替换 $\Phi$ 而忽略 $G_r$。

\subsection{未来温度相关 EM 参数扩展}
若未来引入真正温度相关的海水电导率或介电参数，则 EM surrogate 输入扩展为 $(s_{\rm em},g)$，并在 thermal trajectory 中按需要更新 EM tensors。该扩展必须由 constitutive sensitivity 支持；当前材料模型下不人为增加不存在的热状态依赖。
